% !TeX program = pdflatex

% !TeX root = PF-06.tex

% Inventor (per express statement): Tero Keski-Valkama.

% Review route and source support in REVIEW_NOTES.md before filing.

\documentclass[pct]{sm-patent}

% Change pct to us for US Letter text. A4 drawings are separate.

\PatentTitle{Local learning through event-transported eligibility and independently updated feedback}

\begin{document}

\descriptionstart

\pheading{Technical field}

\pnum Learning circuits and processors having timestamped decaying and rotating state.

\pheading{Background}

\pnum A history-based reverse pass can require stored trajectories and coordinated backward weight access. Replacing exact traces by layer-local approximations changes the derivative when upper memories depend on lower parameters.

\pheading{Summary}

\pnum A temporal learner transports parameter sensitivities with event-state dynamics, includes cross-layer influence in upper sensitivities, and contracts a local state error with those sensitivities. A feedback embodiment maintains independent matrices through coordinated updates.

\pheading{Brief description of the drawings}

\pnum FIG. 1 is a block diagram of an implementation. FIG. 2 shows operations of the implementation. Reference signs identify the same functional elements throughout the drawings.

\pheading{Detailed description}

\pnum In FIG. 1, 100 identifies events; 110 identifies mode state; 120 identifies eligibility; 130 identifies contraction; 140 identifies update; 150 identifies feedback. In FIG. 2, 210 identifies transport, 220 identifies local error and 230 identifies bank update.

\pnum For a mode $z_n=a_nz_{n-1}+b_n$, let $a_n=\exp((-r+\mathrm{i}\omega)\Delta t_n)$. State storage 110 and trace storage 120 retain $z$ and its sensitivity $S=\partial z/\partial\theta$ at the relevant parameter version. For fixed-version trajectory differentiation,
\[S_n=a_n S_{n-1}+(\partial_\theta a_n)z_{n-1}+\partial_\theta b_n.\]
For example $\partial_r a_n=-\Delta t_n a_n$ and $\partial_\omega a_n=\mathrm{i}\Delta t_n a_n$.

\pnum For a local real-valued loss, update engine 130 contracts its real state cotangent with the corresponding real-coordinate sensitivity. Complex notation is implemented as two real coordinates. A reverse embodiment transports a cotangent with the conjugate coefficient $a_n^*$ in reverse event order.

\pnum For stacked memories, an upper trace includes derivatives through lower inputs. If an upper write is $b_2(u)$, its trace includes $(\partial b_2/\partial u)(\partial u/\partial\theta)$ in addition to its own parameter derivatives. Retaining only upper parameters is an explicitly approximate alternative, not exact cross-layer learning.

\pnum A feedback bank can store $B$ independently of a forward matrix $W$. With compatible orientation, both receive the same locally computed increment and the same decoupled shrinkage: $W'=(1-\eta d)W+\Delta$ and $B'=(1-\eta d)B+\Delta$. Their difference contracts by $1-\eta d$ when its magnitude is below one. Equality initially gives equality under these updates; arbitrary nonlinear optimizer differences need separate treatment.

\pnum When a parameter changes during a stream, the processor declares whether traces target the newly evolving learner or a fixed issued computation. It invalidates or transports affected traces according to that declaration. A fixed-weight derivative identity does not by itself justify reusing traces after arbitrary intervening updates.

\pnum FIG. 1 shows event input 100, state 110, sensitivity 120, contraction 130, update output 140 and feedback bank 150. FIG. 2 shows transport 210, local contraction 220 and update 230. A circuit embodiment colocates mode state and sensitivities with multipliers. Trace dimensions and cross-layer products are included in resource counts; history-independent storage can still be large in parameter count.

\pheading{Hardware implementation}

\pnum A digital implementation may use synthesizable register-transfer logic mapped to an FPGA or to an application-specific integrated circuit. Addressed state uses register banks or SRAM macros; arithmetic uses integer, fixed-point or floating-point datapaths. Ready/valid or request/acknowledge interfaces apply backpressure without accepting a message twice. Synchronous FPGA execution, locally clocked circuits and asynchronous circuits are alternative implementations, each preserving its declared causal ordering and numerical precision.

\pnum Before deployment, a reference trace can check selected addresses, values, timestamps, publication versions and output probabilities against the circuit at the declared numerical precision. Timing quantization, ties, cancellation, overflow and saturation are specified for the implementation. FPGA resource utilization, post-layout timing and measured energy are implementation-dependent quantities; no particular measured hardware result is required by this embodiment.

\claimsstart

\begin{patentclaims}

\item A method of operating an event-driven learning processor, comprising maintaining decaying and rotating timestamped memory modes and parameter sensitivities; transporting each sensitivity using a coefficient used for its associated memory mode and derivatives of a state transition; including in an upper-memory sensitivity an influence of a lower-memory parameter through an inter-layer input; contracting a local state error with the sensitivities; and applying a resulting parameter update.

\item The method of claim 1, wherein real and imaginary coordinates implement a complex decay-and-rotation coefficient.

\item The method of claim 1, wherein a reverse cotangent is transported with a conjugate state-transport coefficient.

\item The method of claim 1, further comprising maintaining an independent feedback matrix and applying coordinated local increments and decoupled decay to the feedback matrix and a forward matrix.

\item The method of claim 1, wherein a trace carries an issued parameter version and is invalidated or transported when an influencing parameter changes.

\item A learning apparatus comprising timestamped mode storage, sensitivity storage, a transport datapath and an update engine configured to perform the method of claim 1.

\item A non-transitory computer-readable storage medium storing instructions that cause a processor to perform the method of claim 1.

\end{patentclaims}

\abstractstart

An event-driven learner maintains timestamped memory modes and parameter sensitivities. State transport coefficients also transport sensitivities, including cross-layer influence from lower parameters into upper memories. A local state error is contracted with the sensitivities to obtain an update. A reverse embodiment uses conjugate transport. An independent feedback bank can receive coordinated increments and decay with a forward bank. Version information controls trace reuse after parameter changes. The processor separates history storage savings from the parameter-dependent size of the retained traces.

\end{document}

